\section{枢轴量把未知参数隔离出来}
\label{sec:11-Mathematical-Statistics/04-Interval-Estimation/02}

Wald 区间用渐近正态拼出一个区间。思路直截了当，但边界上的表现已经被上一节说清楚了：\(S=0\) 时区间坍缩成 \([0,0]\)，标称 95\% 成了一个笑话。

有没有办法造出一个覆盖率和标称值精确相等---而不是只能指望大样本---的区间？答案藏在一个巧妙的观察里：如果你能找到一个同时依赖数据和参数的函数，它的分布却是完全已知、干干净净不包含任何未知量的，你就可以先把这个函数的概率事件写下来，然后通过代数整理把参数反解出来。

这样的函数就是枢轴量。

\subsection{公式的反向：从概率事件到区间}

设 \(Q(X,\theta)\) 是一个枢轴量。``枢轴量''的意思是：对每一个可能的 \(\theta\)，\(Q(X,\theta)\) 作为随机变量（随机性来自 \(X\)）的分布是完全已知的，不依赖 \(\theta\) 或任何其他未知参数。

有了枢轴量，构造区间的步骤只有两步。

第一步：利用枢轴量的已知分布，选取分位数 \(a,b\) 使得

\[
P_\theta\{a\le Q(X,\theta)\le b\}=1-\alpha.
\]

第二步：把不等式 \(a\le Q(X,\theta)\le b\) 对 \(\theta\) 反解，得到等价形式

\[
L(X)\le\theta\le U(X).
\]

因为两步之间做的是等价变形，最终事件与最初事件的概率相同：\(P_\theta\{L(X)\le\theta\le U(X)\}=1-\alpha\)。\([L(X),U(X)]\) 就是 \(1-\alpha\) 水平的置信区间。

这个构造的美感在于：你只在一个地方使用了概率---在枢轴量上写下一个关于已知分布的事件。概率计算和参数反解完全分离。你不需要在最后一步对任何量做近似，也不需要祈祷样本量够大。

\subsection{已知方差：最简单的精确区间}

正态数据，\(\sigma\) 已知。

\[
X_i\overset{\mathrm{iid}}{\sim}N(\mu,\sigma^2),\qquad
\bar X\sim N\!\left(\mu,\frac{\sigma^2}{n}\right).
\]

标准化后得到枢轴量：

\[
Z = \frac{\bar X-\mu}{\sigma/\sqrt n}\sim N(0,1).
\]

\(Z\) 的分布是标准正态，与 \(\mu\) 无关。取对称分位数 \(z_{1-\alpha/2}\)：

\[
P\!\left(-z_{1-\alpha/2}\le Z\le z_{1-\alpha/2}\right)=1-\alpha.
\]

现在把不等式展开并对 \(\mu\) 反解。左边：

\[
-z_{1-\alpha/2}\le\frac{\bar X-\mu}{\sigma/\sqrt n}
\;\Longleftrightarrow\;
\mu\le\bar X+z_{1-\alpha/2}\frac{\sigma}{\sqrt n}.
\]

右边：

\[
\frac{\bar X-\mu}{\sigma/\sqrt n}\le z_{1-\alpha/2}
\;\Longleftrightarrow\;
\bar X-z_{1-\alpha/2}\frac{\sigma}{\sqrt n}\le\mu.
\]

两步都是乘正数 \(\sigma/\sqrt n>0\)，不等式方向保持。合在一起：

\[
P_\mu\!\left(
\bar X-z_{1-\alpha/2}\frac{\sigma}{\sqrt n}
\le\mu\le
\bar X+z_{1-\alpha/2}\frac{\sigma}{\sqrt n}
\right)=1-\alpha.
\]

区间为 \(\bar X\pm z_{1-\alpha/2}\,\sigma/\sqrt n\)。整个推导中没有任何近似：正态假设下，覆盖率精确等于 \(1-\alpha\)。

\subsection{未知方差：为什么冒出一个 \(t\) 分布}

\(\sigma\) 未知时，最自然的想法是把 \(\sigma\) 换成样本标准差 \(S\)：

\[
S^2 = \frac{1}{n-1}\sum_i (X_i-\bar X)^2.
\]

问题是 \(\frac{\bar X-\mu}{S/\sqrt n}\) 不再服从正态分布---分母里多了一个随机量 \(S\)。

但这里发生了一件正态模型特有的好事。把分子分母同除以 \(\sigma\)：

\[
\frac{\bar X-\mu}{S/\sqrt n}
= \frac{\frac{\bar X-\mu}{\sigma/\sqrt n}}
{\sqrt{\frac{(n-1)S^2}{\sigma^2}\big/(n-1)}}
= \frac{Z}{\sqrt{V/(n-1)}},
\]

其中

\[
Z=\frac{\bar X-\mu}{\sigma/\sqrt n}\sim N(0,1),\qquad
V=\frac{(n-1)S^2}{\sigma^2}\sim\chi^2_{n-1},
\]

而且 \(Z\) 和 \(V\) 独立。

标准正态除以独立卡方除以其自由度的平方根---这个结构恰好就是 \(t\) 分布的定义。于是

\[
\frac{\bar X-\mu}{S/\sqrt n}\sim t_{n-1},
\]

分布的自由度是 \(n-1\)，与未知的 \(\mu\) 和 \(\sigma\) 都无关。这个量是枢轴量。

有了它，区间立即得到：

\[
\bar X \pm t_{n-1,\,1-\alpha/2}\,\frac{S}{\sqrt n},
\]

其中 \(t_{n-1,\,1-\alpha/2}\) 是 \(t_{n-1}\) 分布的 \(1-\alpha/2\) 分位数。

这里有两个隐藏条件需要摊开来看。第一，\((n-1)S^2/\sigma^2\sim\chi^2_{n-1}\) 依赖正态假设，不只是一个中心极限定理的后果。第二，\(\bar X\) 与 \(S^2\) 的独立性也是正态 iid 模型的特有性质。两个条件合在一起，才给了你一个精确的 \(t\) 分布。``方差未知就用 \(t\)''这句口诀的背后，站着一个完整的正态假设。非正态小样本下，\(t\) 区间只是一个近似，且近似的可靠程度需要另外验证。

\(t\) 分位数比标准正态分位数大，而且自由度越小差得越多。多出来的那部分宽度，正在为你支付``连方差也要从同一批数据里估计''的代价。样本量增大时，\(t_{n-1}\) 趋近 \(N(0,1)\)，这笔额外费用逐渐消失---因为 \(S\) 对 \(\sigma\) 的估计越来越准了。

\subsection{正态方差的卡方区间：不对称来自参数几何}

同一正态模型下，方差 \(\sigma^2\) 本身也有枢轴量：

\[
\frac{(n-1)S^2}{\sigma^2}\sim\chi^2_{n-1}.
\]

取卡方上下分位数 \(\chi^2_{n-1,\,\alpha/2}\) 和 \(\chi^2_{n-1,\,1-\alpha/2}\)：

\[
P\!\left(
\chi^2_{n-1,\,\alpha/2}
\le\frac{(n-1)S^2}{\sigma^2}
\le\chi^2_{n-1,\,1-\alpha/2}
\right)=1-\alpha.
\]

反解 \(\sigma^2\) 时，取倒数会翻转不等号方向：

\[
\frac{1}{\chi^2_{n-1,\,1-\alpha/2}}
\le\frac{\sigma^2}{(n-1)S^2}
\le\frac{1}{\chi^2_{n-1,\,\alpha/2}}.
\]

最终区间：

\[
\left[
\frac{(n-1)S^2}{\chi^2_{n-1,\,1-\alpha/2}},\;
\frac{(n-1)S^2}{\chi^2_{n-1,\,\alpha/2}}
\right].
\]

注意分母的上下分位数位置互换了---这是因为取完倒数后大小关系颠倒。这个区间不对称：下界离 \(S^2\) 比上界近。不是公式写错了，而是方差参数恒为正，卡方分布天然右偏。强行使用``估计值加减若干倍标准误''的对称形式，会生产出延伸到负数的荒谬下界。参数空间的几何决定了哪些区间形状才是合理的。

\subsection{同一件事的两面读法：检验和区间的对偶}

区间构造和假设检验共享同一套骨架。对候选参数 \(\theta_0\) 做一个显著性水平 \(\alpha\) 的检验，把所有``不被拒绝''的 \(\theta_0\) 收集起来，你就得到了一个 \(1-\alpha\) 的置信集合。

正态均值的例子最简单：\(z\) 检验的不拒绝条件是

\[
\left|\frac{\bar X-\mu_0}{\sigma/\sqrt n}\right|\le z_{1-\alpha/2}.
\]

把它当成对 \(\mu_0\) 的约束来解，恰好给出 \(z\) 区间。

这条对偶关系不仅解释了构造上的相似性，还传递了有限样本行为：Wald 检验的毛病（边界处实际 size 偏离标称水平）会原封不动地传染给 Wald 区间；score 检验的好行为同样会让 score 区间更靠谱。选择区间方法时，等于在选择背后那条检验的逻辑。

\subsection{把不等式方向盯紧}

从概率事件到区间端点的代数反解，最大的陷阱是方向。每一步乘法或取倒数，都需要检查乘数或除数的符号：

\begin{itemize}
\tightlist
\item 乘正数：不等号方向不变；
\item 乘负数：不等号翻转；
\item 取倒数（两边同号）：不等号翻转；
\item 取倒数（两边异号或含零）：事件必须分段处理。
\end{itemize}

方差卡方区间中下分位数出现在上界、上分位数出现在下界，就是因为取倒数的翻转效应。最可靠的办法是先写下完整的双侧事件，再一步一步做等价变形，每步确认方向。背诵公式---尤其是不对称区间的公式---很难避免在某次使用时把分位数放错位置。

\subsection{精密不等于准确：秤的故事}

同一个行李箱，你每天放上家里的电子秤连称 10 次。读数极其一致，\(S\) 很小，\(t\) 区间很窄。你获得了一个高度精密的均值区间。

但如果这把秤始终多报 \(0.4\,\mathrm{kg}\) 的系统偏差呢？十次读数的集中性只缩小了随机项 \(S/\sqrt n\)，和多出来的那 \(0.4\,\mathrm{kg}\) 毫无关系。增加重复次数不能消掉共同偏差。你的 \(t\) 区间可以非常窄，却可以系统地错过真实重量。

实验室传感器面对同样的困境。要评价准确度，不能只靠同一台设备在同一条件下的重复读数。还需要可追溯的参考标准、跨时间和环境的校准设计，或把漂移和批次效应直接写成模型的一部分。覆盖率来自完整的测量机制---从数据采集到模型设定到算法选择---而不来自区间端点的那个公式本身。

枢轴量方法给了一条不需要渐近凑合的精确路径。但它要求你真正理解枢轴量成立的条件：什么分布假设在支撑那个已知分布，不等式方向每次是怎么变的，以及重复测量的精密是否真的对应着你关心的那个量。这些检查做完之后，你算出来的区间，才配得上那个``置信''二字。
